\lesson{project}{先读懂项目，再决定怎么改}
\diagram{p_project}
\topic{先调研已有方案，再讨论架构}
打开正确项目，先查已有实现和相似项目，再决定复用、改造或自建。依据官方文档、原始研究和示例，比较适用场景、成本及限制。讨论模块分工、数据流与接口（程序间交换数据的约定）；追问“为什么这样设计、何时不适用”，让新想法建立在已有工作之上。
\topic{开工前保护好数据}
\begin{enumerate}
\item \textbf{用测试环境。}用数据副本和测试账号，限制读写与网络访问。扩大权限试验时隔离正式系统；删除、发布或修改正式数据前，明确授权、备份和恢复办法。
\item \textbf{保管密钥。}优先用密钥管理服务；本地 \texttt{.env} 仍需限制访问且不得提交，不贴进聊天或网页代码。\texttt{.gitignore} 不影响已跟踪文件，不阻止 AI 读取，也不清除已有历史；泄露后撤销并更换密钥。
\item \textbf{先查陌生内容。}网页或文件中的指令诱导 AI 偏离任务，叫提示注入。读取时避免同时开放私密数据与外发、删除权限；第三方技能、脚本和软件包先核对来源。
\end{enumerate}
\begin{prompt}{项目阅读与改动定位}
目标：[当前任务]，材料：[位置]，允许更新的调研文件：[路径]。核对版本、已有改动、规则和运行条件；本轮不改程序，只在测试环境验证。材料不可读先报缺口，不能写记录则返回草稿。\par
查找已有实现与相似方案，比较复用、改造与自建；核对软件包、官方来源和使用许可。无法联网时标明外部资料未核实，不编造来源。围绕架构未知做最小验证，说明模块、数据流与依赖。\par
记录问题、来源与日期、观察、备选、选择理由及未知，区分事实、推断和假设；说明什么新证据会改变决定。报告实际改动与未验证项，保留原始资料和检查证据。
\end{prompt}
\principle{让调研形成决定}{上下文是 AI 当前能看到的信息，应按需提供。调研要回答“能复用什么、为什么选它、什么还需验证”，并落实到方案。文件存在不等于已经读到；本地运行也不代表资料不会传给云端模型。}
